\documentclass{article}
\usepackage{amsmath, amssymb, amsthm}

\title{IMC 2021, Second Day Solutions}
\date{August 4, 2021}

\begin{document}
\maketitle

\section*{Problem 5}
\subsection*{Solution}
Let $B_m$ be the matrix with $2021 B_m = A^m + B_m^2$. For $m = 1$: $A = 2021 B_1 - B_1^2$, so $A$ is real symmetric, hence diagonalizable with real eigenvalues.

For any eigenvalue $\lambda$ of $A$ and corresponding eigenvalue $\mu$ of $B_m$: $\mu^2 - 2021\mu + \lambda^m = 0$. Real discriminant: $2021^2 - 4\lambda^m \geq 0$, so $\lambda^m \leq 2021^2/4$. If $|\lambda| > 1$, taking $m$ even and large gives contradiction. Hence $|\lambda| \leq 1$ for all eigenvalues, so $|\det A| \leq 1$.

\section*{Problem 6}
\subsection*{Solution}
For $p = 2$: $|\mathrm{GL}_2(\mathbb{Z}/2)| = 6 > 2 = |S_2|$.

For odd $p$: $A = \begin{pmatrix} 1 & 1 \\ 0 & 1 \end{pmatrix}$ has order $p$, and commutes with $B = -I$ of order $2$. So $AB$ has order $2p$. But in $S_p$, any element of order divisible by $p$ must be a $p$-cycle (order exactly $p$), so no element has order $2p$. Contradiction.

\section*{Problem 7}
\subsection*{Solution}
Let $p(z) = z^n + a_{n-1} z^{n-1} + \cdots + a_0$ and $q(z) = z^n \overline{p(1/\bar z)} = 1 + \overline{a_{n-1}} z + \cdots + \overline{a_0} z^n$. For $|z| = 1$, $1/\bar z = z$, so $|q(z)| = |p(z)|$. Also $q(0) = 1$. By maximum modulus applied to $f \cdot q$:
\[ |f(0)| = |f(0) q(0)| \leq \max_{|z|=1} |f(z) q(z)| = \max_{|z|=1} |f(z) p(z)|. \]

\section*{Problem 8}
\subsection*{Solution}
Answer: $2n$.

\textbf{Example.} A basis $e_1, \ldots, e_n$ and $-e_1, \ldots, -e_n$: from any 3, at least 2 are orthogonal (only pair $e_i, -e_i$ is non-orthogonal).

\textbf{Upper bound (Solution 2).} Let $P_i$ be the rank-1 projection onto the $i$-th vector. The hypothesis gives $\operatorname{tr}(P_i P_j P_k) = 0$ for distinct $i, j, k$. Let $Q = \sum P_i$, with eigenvalues $t_1, \ldots, t_n \geq 0$ and $\sum t_i = \operatorname{tr}(Q) = N$.
\[ \sum t_i^3 = \operatorname{tr}(Q^3) = N + 6 \sum_{i<j} \operatorname{tr}(P_i P_j) = N + 3(\operatorname{tr}(Q^2) - N) = 3 \sum t_i^2 - 2N. \]
Since $(t-2)^2(t+1) \geq 0$ gives $t^3 - 3t^2 + 4 \geq 0$:
\[ -2N = \sum (t_i^3 - 3 t_i^2) \geq -4n, \]
so $N \leq 2n$.

\end{document}
